%% kapa: exam=1 round=28 year=2017 period=1 form=A subject=economics subject_ko=경제학원론 seq=36 number=76
%% source: Q-Net 감정평가사 기출문제 (cst00309) artlSeq=5199603 file=extracted/1008/5199603/2017년 제28회 감정평가사 1차시험 문제지/1교시 A형.pdf
%% conversion: kordoc 4.24.51 -> tools/gen1.py (auto); review=proofread-sonnet; flags=none
\begin{question}{76}
총수요-총공급 모형의 단기 균형 분석에 관한 설명으로 옳은 것은? (단, 총수요곡선은 우하향하고, 총공급곡선은 우상향한다.)
\begin{choices}
  \item 물가수준이 하락하면 총수요곡선이 오른쪽으로 이동하여 총생산은 증가된다.
  \item 단기적인 경기변동이 총수요충격으로 발생되면 물가수준은 경기역행적(counter-cyclical)으로 변동한다.
  \item 정부지출이 증가하면 총공급곡선이 오른쪽으로 이동하여 총생산은 증가한다.
  \item 에너지가격의 상승과 같은 음(-)의 공급충격은 총공급곡선을 오른쪽으로 이동시켜 총생산은 감소된다.
  \item 중앙은행이 민간 보유 국채를 대량 매입하면 총수요곡선이 오른쪽으로 이동하여 총생산은 증가한다.
\end{choices}
\end{question}
\answer{5}
